\section{Simple linear regression}
Now we switch focus on \textit{supervised learning}: that is, deriving/predicting the value of some feature given the values of others.
Taking this example on \textcolor{green}{exam score vs study hours}, what if we'd like to predict the score
of the next person who studies 5 hours per week?
\begin{figure}[h]
    \centering
    \includegraphics[width=0.3\textwidth]{linearmodel/lm_example.png}
\end{figure}

Obviously we take as predcitor the mean of the score samples for that amount of studying hours; we can do the same for the other sets of points. We can see there is a linear model between the 
study hours and the mean of the exam scores: $\mathbb{E}[\text{score}|\text{study\ hours}] = f(\text{study\ hours})$.
\vspace{0.2cm}

So we want to create a model to shape this relation:
\begin{description}
    \item[] $X$ := independent/predictor variable (usually known)
    \item[] $Y$ := dependent/response variable \arrow it's a random variable (due \seps)
    \item[] \seps := random deviation/error \arrow it's a random variable
\end{description}
\begin{equation*}
    \mathbb{E}[Y|X]=\beta_0+\beta_1X+\varepsilon
\end{equation*}
where $\beta_0,\beta_1,\sigma^2$ are parameters and we \textbf{assume} that $\mathbb{\varepsilon}=0$
,$\mathbb{V[\varepsilon]}=\sigma^2$ and that's symmetrically distributed.\\ 
\vspace{0.2cm}
Interpretation of the parameters:
\begin{itemize}
    \item \textbf{intercept} $\beta_0$: "height"/base-line of the model, where the line intercepts the 0
    \item \textbf{slope} $\beta_1$: how much my model increases for unit increment of the 
    predictor $x$ \textcolor{gray}{(by studying 1 hour more, I will improve my score by $beta_1$)}
    \item $\sigma^2$: how much our response is spread along the model prediction
\end{itemize}
\begin{figure}[h]
    \centering
    \includegraphics[width=0.4\textwidth]{linearmodel/lm_param_interpretation.png}
\end{figure}
\newpage
Let's start with the following assumptions on the errors: we assume that the spread of our model on each predictor value is the same (\textit{homoscedasticity}) and that $\varepsilon \sim \mathcal{N}(0,\sigma^2)$.
\begin{figure}[h]
    \centering
    \includegraphics[width=0.4\textwidth]{linearmodel/err_assumption.png}
\end{figure}

\subsection{Model estimation}
The model parameters are not usually known, so how do we estimate them from a set of samples $(x_1,y_1
\dots (x_n,y_n))$?\\
We assume $y_i = \beta_0 +\beta_1 x_i +\varepsilon_i,\ \ \varepsilon_i\sim\mathcal{N}(0,\sigma^2)
\rightarrow y_i\ |\ x_i\sim\mathcal{N}(\beta_0+\beta_1 x_i, \sigma^2)$ and we are going to select those parameters that minimize the squared error:
\begin{equation*}
    \hat{\beta}_0,\hat{\beta}_1 = \arg\min_{\beta}\sum_i\varepsilon_i^2=\sum_i\left(y_i-(\beta_0+
    \beta_1 x_i)^2\right)
\end{equation*}
\begin{figure}[h]
    \centering
    \includegraphics[width=0.4\textwidth]{linearmodel/square_error_visualized.png}
\end{figure}


\textcolor{gray}{(you can skip the computation of the closed formula)}
$$
\begin{cases}
    \frac{\partial}{\partial \beta_0} \sum_{i} (y_i - [\beta_0 + \beta_1 x_i])^2 = -2 \sum_{i} (y_i - \beta_0 - \beta_1 x_i) = 0 \\
    \frac{\partial}{\partial \beta_1} \sum_{i} (y_i - [\beta_0 + \beta_1 x_i])^2 = -2 \sum_{i} x_i (y_i - \beta_0 - \beta_1 x_i) = 0
\end{cases} \implies
$$ $$
\begin{cases}
    \sum_{i} y_i = n\beta_0 + \beta_1 \sum_{i} x_i \\
    \sum_{i} y_i x_i - \beta_0 \sum_{i} x_i - \beta_1 \sum_{i} x_i^2 = 0
\end{cases} \implies
$$ $$
\begin{cases}
\bar{y} = \beta_0 + \beta_1 \bar{x} \\
\sum_{i} x_i^2 \beta_1 = \sum_{i} y_i x_i - (\bar{y} - \beta_1 \bar{x}) \sum_{i} x_i
\end{cases} \implies
$$ $$
\begin{cases}
    \bar{y} = \beta_0 + \beta_1 \bar{x} \\
    \beta_1 \left( \sum_{i} x_i^2 - \frac{(\sum_{i} x_i)^2}{n} \right) = \sum_{i} x_i y_i - \frac{\sum_{i} x_i \sum_{i} y_i}{n}
\end{cases} \implies
$$ $$
\begin{cases}
    \hat{\beta}_0 = \bar{y} - \hat{\beta}_1 \bar{x} \\
\hat{\beta}_1 = \cfrac{\sum_{i=1}^n (x_i - \bar{x})(y_i - \bar{y})}{\sum_{i=1}^n (x_i - \bar{x})^2}
\end{cases}
$$

Even if it might look like PCA, now we're not minimizing the distance of the line from the point
\textcolor{gray}{(is not even a distance, is not orthogonal)} but we're just computing the error 
between the model and the actual data. 

\subsection{Model assessment}
\textbf{After estimating our model, we'd like to check its goodness}.\\
To do so, first we compute the variance of out estimations $\hat{\beta}_0,\hat{\beta}_1$ and then we use it to build a \textbf{confidence interval CI}.

$$
\begin{cases}
Var(\hat{\beta}_0) = \sigma^2 \left[ \frac{1}{n} + \cfrac{\bar{x}^2}{\sum_{i=1}^n (x_i - \bar{x})^2} \right] \\
Var(\hat{\beta}_1) = \cfrac{\sigma^2}{\sum_{i=1}^n (x_i - \bar{x})^2}
\end{cases} \implies
\begin{cases}
\hat{\beta}_0 \sim \mathcal{N}(\beta_0, Var(\hat{\beta}_0)) \\
\hat{\beta}_1 \sim \mathcal{N}(\beta_1, Var(\hat{\beta}_1))
\end{cases}
$$
with an $\alpha$\% CI $\hat{\beta_i}\pm t_{n-2,1-\alpha/2}\sqrt{Var(\hat{\beta}_i)}$

\vspace{0.2cm}
One critical aspect to check is whether the computed CI contains zero or not: if it does the coefficient might be 0, meaning that $Y$ and $X$ are not really dependent \arrow if the CI is completely above/below zero, even by a small amount, no problem.
\vspace{0.1cm}
\begin{figure}[h]
    \centering
    \includegraphics[width=1\linewidth]{linearmodel/lm_param_CI.png}
\end{figure}

\vspace{0.4cm}
We could also approach the problem from another perspective using \textbf{tests}: we'd like to quantify how much we are sure that the estimate won't be zero.
\begin{equation*}
    H_0: \beta_i=0\ \ \text{vs}\ \   H_1: \beta_i\ne0
\end{equation*}

and the statistic used for the test (below) measures the dispersion of the coefficient estimate. If we get small t p-value under the hypo that the coefficient is zero, it means that I'm seeing
something very unlikely\arrow I have evidence that the coefficient is not 0. 
 
\begin{equation*}
    t_i = \cfrac{\hat{\beta_i}}{SE(\hat{\beta_i}) = \sqrt{Var(\hat{\beta_i})}}\sim t_{n-2}
\end{equation*}

\newpage
\subsubsection{Variability explained}
After the parameter estimation, we can compute the \textbf{residuals} $\hat{e}$ := difference 
between our model and the real output.
$$
\begin{aligned}
\hat{y}_i &\to \hat{\beta}_0 + \hat{\beta}_1 x_i + e_i \\
y_i &\to \beta_0 + \beta_1 x_i + \epsilon_i
\end{aligned}
\implies e_i \neq \epsilon_i
$$

\vspace{0.2cm}
\begin{wrapfigure}[14]{l}{0.5\linewidth}
    \includegraphics[width=0.4\textwidth]{linearmodel/var_explained.png}
\end{wrapfigure}
If we don't know any model, the only reasonable way to estimate $Y$ is to use $\bar{y}$.
\vspace{0.2cm}
$$
\begin{aligned}
    &\sum_{i=1}^{n} \textcolor{orange}{(y_i - \bar{y})^2} = 
    \sum_{i=1}^{n} \textcolor{green}{(\hat{y}_i - \bar{y})^2} + 
    \sum_{i=1}^{n} \textcolor{blue}{(y_i - \hat{y}_i)^2} \\
    &\begin{matrix}
        \textcolor{orange}{\textbf{TSS}} \\
        \textcolor{orange}{\text{Total Sum of}} \\
        \textcolor{orange}{\text{Squares}}
    \end{matrix} 
    \ =\ 
    \begin{matrix}
        \textcolor{green}{\textbf{ESS}} \\
        \textcolor{green}{\text{Explained Sum}} \\
        \textcolor{green}{\text{of Squares}}
    \end{matrix}
    \ +\ 
    \begin{matrix}
        \textcolor{blue}{\mathbf{RSS}} \\
        \textcolor{blue}{\text{Residual Sum of}} \\
        \textcolor{blue}{\text{Squares}}
    \end{matrix}
\end{aligned}
$$
\vspace{0.3cm}
\begin{description}
    \item[\textcolor{orange}{TSS}]\arrow total variability of the data, is the biggest mistake we
    can make during prediction
    \item[\textcolor{green}{ESS}]\arrow improvement of the model with respect to the simplest predictor ($\bar{y}$), is the amount of data variability captured by our model
    \item[\textcolor{blue}{RSS}]\arrow difference between our model and the real data (error model),
    is the amount of variability we couldn't capture
\end{description}

\vspace{1.8cm}
So to evaluate the goodness of our model we could measure "the \% of variability we capture with our model":
\begin{equation*}
    R^2 = \frac{\textcolor{green}{ESS}}{\textcolor{orange}{TSS}} = 
    1 - \frac{\textcolor{blue}{RSS}}{\textcolor{orange}{TSS}}
\end{equation*}
We can also used an \textit{adjusted version} that penalizes more complex models, making the adoption of a more complex model not always the best option:
\begin{equation*}
    R^2_{adj} = 1-\cfrac{\textcolor{blue}{RRS}/(n-p)}{\textcolor{orange}{TSS}/(n-1)}
\end{equation*}


\subsubsection{Residual variability/prediction error}
We would like to assess a also the variability of our model prediction. To do so, there are several options:
\begin{description}
    \item[Mean Absolute Error (MAE)] 
    plain error average.
    \begin{equation*}
        MAE = \frac{1}{n}\sum_{i=1}^{n}|y_i-\hat{y}_i|
    \end{equation*}
    
    \item[Root mean squared error (RMSE)]
    average prediction error with higher emphasis on large mistakes.    
    \begin{equation*}
        RMSE = \sqrt{\frac{1}{n}RSS}
    \end{equation*}

    \item[Residual Standard Error (RSE)] 
    similar to RMSE, but now we have a penalization for more complex models.
    \begin{equation*}
        \hat{\sigma}^2=\cfrac{RSS}{n-p}
    \end{equation*}
    

\end{description}


\subsubsection{Model usefulness}
Another thing we'd like to evalate is how useful our model is in explaining the relationship between the response and the predictor: our model is useful if $\beta_1\ne 0$, and the more variability is able to caputure the better \arrow we use an \textbf{F-test}.
\begin{equation*}
    H_0:\beta_1 = 0 \rightarrow F=\cfrac{ESS/(p-1)}{RSS/(n-p)}\sim F_{(p-1)(n-p)}
\end{equation*}
Trying to give it an intepretation, if we consider:
\begin{itemize}
    \item $ESS/(p-1)$ as the average explained variability per each parameter
    \item $RSS/(n-p)$ as the average unexplained variability per residual degree of freedom of the data
\end{itemize}
we get that \textit{$F$ compares the "information per parameter" with the "noise per degree of freedom"}.\\
If the statistic is low then there are evidence against the null hypo. Note that in the linear model, the F-test $\equiv$ t-test, giving the same result.


\subsection{Inference on Y}
Our model coefficient are only an estimation of the true system's parameters (the fitted lines has some uncertanty) so instead of point-level prediction we'd like to use intervals to predict our response.

\subsubsection{CI for the mean response}
The problem we're going to answer is in the form "what's the avg exam score of students who studies $x_0$ hours?" \arrow we want to estimate $\mu_{Y|X=x_0}=\mathbb{E}[Y|X=x_0]$ and to do so we use CI.

\begin{equation*}
    \hat{y}_0 \pm t_{n-2,1-\alpha/2}\ \hat{\sigma}\sqrt{\frac{1}{n}+\cfrac{(x_0-\bar{x})^2}
    {\sum_{i=1}^n(x_i-\bar{x})^2}}
\end{equation*}
\begin{figure}[h]
    \centering
    \includegraphics[width=0.4\textwidth]{linearmodel/CI_mean.png}
\end{figure}

We could see that the interval is more narrow near $x_0=\bar{x}$, and makes sense: the further we shift from our dataset center of mass the less density of data, incrementing therefore our uncertanty in those regions.

\subsubsection{Prediction interval}
Now we want to tackle the problem "if the next student studies $x_0$ hours, what range of scores it could hope for?"\arrow we want to estimate $Y_{new}|X=x_0$. In this scenario we have the uncertanty due our model parameters (as before, they're only ans estimation) and due the irreducible noise around the reponse mean \seps.
\begin{equation*}
    \hat{y}_0 \pm t_{n-2,1-\alpha/2}\ \hat{\sigma}\sqrt{1+\frac{1}{n}+\cfrac{(x_0-\bar{x})^2}
    {\sum_{i=1}^n(x_i-\bar{x})^2}}
\end{equation*}


\subsubsection{Comparison: CI vs PI}
\begin{description}
    \item[CIs] show the uncertanty of out model on the population/true system parameter, and the interval becomes larger as the sample's std deviation does (so this issue could be reduced by using a larger sampleset)
    \item[PIs] show the uncertanty on some \textit{specific} value, not on the entire dataset, thus the resulting interval is much wider compared to the CIs
\end{description}

\begin{figure}[h]
    \centering
    \includegraphics[width=0.75\textwidth]{linearmodel/lm_CI_vs_PI.png}
\end{figure}



